\chapter{Values, types and variables}\label{ch-values}

This chapter introduces the building blocks of every program: values and
their types, constants and variables, the arithmetic and logical
operators, assignment, and simple input and output. One idea in it is
more important than the others and is the key to much of the language:
a variable in Algol 68 is a name for a \emph{reference} to a memory cell,
and the language says so openly. Section~\ref{sec-references} explains
what this means.

\section{Types}\label{sec-types}

Every value in a program has a \term{type} (\rr{mode}). The type
determines which values are possible and which operations apply. Algol 68
has four \term{primitive types}:

\begin{center}
\begin{tabular}{lll}
\toprule
Type & Values & Examples \\
\midrule
\code|int| & whole numbers & \code|0|, \code|42| \\
\code|real| & floating-point numbers & \code|3.5|, \code|1e6| \\
\code|bool| & truth values & \code|true|, \code|false| \\
\code|char| & single characters & \code|"a"|, \code|"7"| \\
\bottomrule
\end{tabular}
\end{center}

The type \code|string| for sequences of characters is not primitive: it is
a resizable array of \code|char| (Chapter~\ref{ch-advanced}). It can be used in the same
simple ways as the primitive types, though, and this chapter does so.

The range of \code|int| and the precision of \code|real| depend on the
implementation. The standard library provides constants that describe
them, among them \code|max_int|, the largest \code|int|, and
\code|small_real|, the smallest \code|real| that gives a result different
from 1 both when it is added to 1 and when it is subtracted from 1:

\example{ch02/limits}

\begin{implnote}
In ga68 and a68j, \code|int| has 32 bits and \code|real| is a 32-bit
floating-point number (single precision, about 7 decimal digits). The
types \code|long int| and \code|long real| have 64 bits. In Algol 68 Genie,
\code|int| has 64 bits and \code|real| is a 64-bit floating-point number.
Precision variants are the subject of Chapter~\ref{ch-moretypes}.
\end{implnote}

\section{Literals}\label{sec-literals}

A value written directly in the program text is a \term{literal}
(\rr{denotation}):

\example{ch02/literals}

An \code|int| literal is a sequence of digits. It has no sign: in
\code|-5|, the minus is an operator applied to the literal \code|5|.
Blanks may appear between the digits, so \code|1 000 000| is a single
literal; commas may not.

A \code|real| literal has a decimal point, an exponent, or both:
\code|3.5|, \code|.25|, \code|1e6|, \code|2.5e-3|. There must be at least
one digit after the point, so \code|5.| is not a literal; \code|5.0| is.
The letter of the exponent may be written \code|e| or \code|E|.

A \code|char| literal is a single character between double quotes. A
double quote itself is written twice: \code|""""| is the literal for the
character \code|"|. The same rule applies to strings: \code|"say ""hi"""|
is the string \texttt{say "hi"}.

The two \code|bool| literals are \code|true| and \code|false|. When
\code|print| writes a truth value, it writes \code|T| or \code|F|.

\section{Names}\label{sec-names}

The names that a program gives to its constants, variables and procedures
are \term{identifiers} (\rr{identifier}, \rr{tag}). An identifier begins
with a lower-case letter and continues with lower-case letters, digits
and underscores: \code|x|, \code|x2|, \code|max_value|. Reserved words
cannot be used as identifiers.

Capital letters have a meaning of their own (Section~\ref{sec-stropping}): a word that
begins with a capital letter is a type name or an operator. A name in the
style of other languages, such as \code|maxValue|, is therefore read as
the identifier \code|max| followed by a type name \code|Value|, and the
compiler complains about the type:

\example{ch02/identifier-case}

\section{Constants}\label{sec-constants}

A declaration of the form
\begin{lstlisting}
int days_per_week = 7
\end{lstlisting}
binds the identifier \code|days_per_week| to the value 7, for the rest of
the block in which the declaration appears. Such a binding cannot be
changed; \code|days_per_week| is a \term{constant}. The Report calls this
form an \rr{identity declaration}, since it declares that the identifier
and the value are identical: wherever \code|days_per_week| is used, it
stands for 7. It corresponds to \code|const| in C and JavaScript, to
\code|final| in Java and to \code|let| in Rust.

The value need not be known when the program is compiled. Any expression
of the right type can stand on the right-hand side, including one that
uses other constants or reads input. Several constants can be declared
at once, separated by commas; declarations are separated by semicolons:

\example{ch02/constants}

An attempt to change a constant is an error, found by the compiler:

\example{ch02/constant-assign}

The message is phrased in the terms of the Report: a \rr{name} is what
this book calls a reference, and the message says that \code|n| is not a
reference, so nothing can be assigned to it. Section~\ref{sec-references}
explains this.

\section{Variables}\label{sec-variables}

A \term{variable} is declared with \code|:=| instead of \code|=|, or with
no initial value at all:

\example{ch02/variables}

The declaration \code|int count := 0| creates a variable that holds an
\code|int|, with the initial value 0. The \term{assignment}
\code|count := count + 1| replaces the value held by \code|count| with a
new value. Several declarations can be joined with commas, but the type
may be left out only if the next declaration is of the same kind:
\code|int x = 5, y := 6| is an error (``mixed identity-declaration and
variable-declaration''), whereas \code|int x = 5, int y := 6| is allowed.
The declarations joined by commas are carried out in an unspecified
order, so one of them must not use another
(Section~\ref{sec-collateral}).

A variable declared without an initial value, as in \code|real x|, holds
no defined value until a value is assigned to it. Using it before then is
an error that the compiler does not detect. A \code|string| variable, and a \code|string| inside an array or record
variable, is the exception: it starts out as the empty string, since a string is a
resizable array whose initial bounds are 1 and 0
(Section~\ref{sec-flex}), and an empty array has no elements that could
be undefined.

\begin{implnote}
ga68 and a68j give such a variable an initial value: zero, \code|false|
or a blank character. A program should not rely on this, since other
implementations do not.
\end{implnote}

\section{What a variable is: references}\label{sec-references}

In most languages a variable is a primitive concept: a name that refers
to a place in memory whose content can change. Algol 68 builds variables
from two simpler concepts, and makes both visible.

The first concept is \term{allocation}. The expression \code|loc int|
allocates a fresh memory cell that can hold an \code|int|, on the stack
of the current block, and yields a \term{reference} to that cell. The
type of a reference to an \code|int| cell is \code|ref int|, read
``reference to int''. A reference is a value like any other: it can be
bound to an identifier, passed to a procedure or stored in another cell.
The Report calls a reference a \rr{name} and an allocation a
\rr{generator}.

The second concept is the constant declaration of the previous section.
Together they give a variable:

\example{ch02/references}

The declaration \code|ref int b = loc int := 5| allocates a cell, assigns
5 to it, and binds \code|b| to the reference. The short form
\code|int a := 5| means exactly the same, and a program normally uses the
short form. But the long form shows what the identifier of a variable
stands for: not the number in the cell, but the reference to the cell.
The type of \code|a| is \code|ref int|, not \code|int|.

Since a reference is a value, it can be bound to a second name. The
second name then refers to the same cell, and an assignment through it
changes the variable:

\example{ch02/alias}

\noindent prints 99. Nothing has been copied: the declaration binds the
name \code|alias| to the reference that \code|a| stands for.

This explains the error message of the previous section. The left-hand
side of an assignment must be a reference, of type \code|ref int| when an
\code|int| is assigned. A constant declared as \code|int n = 10| has the
type \code|int|; it is the number 10 itself, and there is nothing to
assign to.

Programmers who know C++ will recognize the idea: \code|b| behaves like
a C++ reference \texttt{int\& b} to a cell. In Rust, \code|b| would be a
\texttt{\&mut i32}, which needs an explicit \texttt{*} to read the cell;
Algol 68 reads the cell implicitly where a value is needed
(Section~\ref{sec-coercions}). Programmers who know Java or Python can think of the
reference as the hidden storage location behind every variable; Algol 68
simply gives it a type. The distinction between a value and a reference
to a cell holding that value is the source of most of the power of the
language's type system, and later chapters return to it often: references
to references, references in records, references returned by functions,
and how long the cell behind a reference lives.

\section{Arithmetic}\label{sec-arithmetic}

The arithmetic operators on \code|int| and \code|real| are:

\begin{center}
\begin{tabular}{lll}
\toprule
Operator & Meaning & Operand types \\
\midrule
\code|+| \code|-| \code|*| & addition, subtraction, multiplication & \code|int|, \code|real| \\
\code|/| & division & \code|int|, \code|real|; the result is \code|real| \\
\texttt{\%} & integer division, truncating toward zero & \code|int| \\
\code|MOD| & remainder, never negative & \code|int| \\
\code|**| & power & \code|int| or \code|real| base, \code|int| exponent \\
\code|-| (unary) & negation & \code|int|, \code|real| \\
\bottomrule
\end{tabular}
\end{center}

Note that \code|/| always yields a \code|real|, even for two \code|int|
operands, so \code|n := 4 / 2| is rejected for an \code|int| variable
\code|n|; integer division is \code|%|. The remainder \code|a MOD b| lies
between 0 and \code|ABS b - 1| even if \code|a| is negative, unlike the
\texttt{\%} operator of C, Java or JavaScript:

\example{ch02/arithmetic}

The power operator requires an \code|int| exponent, which may be
negative if the base is a \code|real|: \code|2.0 ** -1| is 0.5. A
\code|real| power,
such as a square root, is computed with the functions of the standard
library (\code|sqrt|, \code|exp|, \code|ln|), which Chapter~\ref{ch-moretypes} describes.

\begin{implnote}
An \code|int| result that does not fit into an \code|int| is undefined by
the Report. In ga68 and a68j, it wraps around silently:
\code|max_int + 1| gives $-2147483648$. The result of an \code|int| raised to a
negative power, such as \code|2 ** -1|, is undefined as well.
\end{implnote}

\section{Precedence}\label{sec-precedence}

Every binary operator has a \term{precedence} (\rr{priority}) from 1 to 9.
An operator of higher precedence binds more tightly: \code|2 + 3 * 4| is
\code|2 + (3 * 4)|. Operators of equal precedence are applied from left to
right: \code|10 - 4 - 3| is \code|(10 - 4) - 3|. The standard operators
have the following precedences:

\begin{center}
\begin{tabular}{cl}
\toprule
Precedence & Operators \\
\midrule
9 & \code|I| (complex numbers, Chapter~\ref{ch-moretypes}) \\
8 & \code|**|, \code|UP|, \code|LWB|, \code|UPB| \\
7 & \code|*|, \code|/|, \texttt{\%}, \code|MOD|, \code|ELEM| \\
6 & \code|+|, \code|-| \\
5 & \code|<|, \code|<=|, \code|>=|, \code|>| \\
4 & \code|=|, \code|/=| \\
3 & \code|AND| \\
2 & \code|OR| \\
1 & \code|+:=|, \code|-:=|, \code|*:=|, \code|/:=|, \texttt{\%:=}, \ldots \\
\bottomrule
\end{tabular}
\end{center}

Two rules differ from what most current languages do. First, \code|**|
is applied from left to right like every other operator, so
\code|2 ** 3 ** 2| is $(2^3)^2 = 64$, not $2^9 = 512$. Second, every
unary operator binds more tightly than every binary operator. In
\code|-2 ** 2|, the minus applies to the 2 first, so the result is
$(-2)^2 = 4$:

\example{ch02/precedence}

When in doubt, parentheses make the intended order explicit, and the
programs in this book use them wherever the order might surprise a reader.
Older programs, and many on Rosetta Code, write some operators as words:
\code|OVER| for \texttt{\%}, \code|PLUSAB| for \code|+:=|, and so on
(Appendix~\ref{app-library}).
Chapter~\ref{ch-procs} shows how a program declares operators of its own and their
precedences.

\section{Comparisons and logical operators}\label{sec-logic}

The comparison operators \code|=|, \code|/=| (not equal), \code|<|,
\code|<=|, \code|>|, \code|>=| compare two numbers, or two characters in
the order of their character codes, and yield a \code|bool|. The logical
operators are \code|AND|, \code|OR| and the unary \code|NOT|:

\example{ch02/logic}

Since \code|AND| has a lower precedence than the comparisons,
\code|x > 0 AND x < 10| needs no parentheses. Applied to two truth
values, \code|=| tests whether they are equal and \code|/=| is the
exclusive or; \code|(a < 0) = (b < 0)| tests whether \code|a| and
\code|b| have the same sign.

\begin{pitfall}
\code|NOT| is a unary operator and binds more tightly than every binary
operator, so \code|NOT x = 7| means \code|(NOT x) = 7|, which the
compiler rejects for an \code|int| \code|x|:
\lstinputlisting[style=output]{examples/ch02/not-prec.out}
\noindent The test is written \code|NOT (x = 7)| or \code|x /= 7|.
Programmers who know Python, where \code|not| binds loosely, fall into
this easily. Similarly, \code|x < y < z| is not a chained comparison:
\code|x < y| is a \code|bool|, which cannot be compared with a number.
\end{pitfall}

\code|AND| and \code|OR| are ordinary operators: both operands are
evaluated before the operator is applied, as with every operator. They do
not ``short-circuit'' like \texttt{\&\&} and \texttt{||} in C or Java. The
GNU dialect of Algol 68, which ga68 and a68j implement, adds the
short-circuit forms \code|andth| and \code|orel|, which evaluate their
right operand only when it is needed. They bind more weakly than every
operator, and a chain of them needs parentheses:
\code|a andth (b andth c)| (Appendix~\ref{app-syntax}). The following program shows the
difference; it uses a small procedure \code|check| that prints a line
before it returns its argument (procedures are the subject of
Chapter~\ref{ch-procs}):

\example{ch02/and-evaluates-both}

\section{Implicit conversions}\label{sec-coercions}

In the program of Section~\ref{sec-references}, \code|a| has the type
\code|ref int|, and yet \code|a + 1| is allowed, although \code|+| is
defined for \code|int| and not for references. The compiler inserts an
operation that reads the cell: it converts the reference into the value it
refers to. Such an automatic change of type is an \term{implicit
conversion} (\rr{coercion}), and this one is called
\term{dereferencing} (\rr{dereferencing}).

Algol 68 has six kinds of implicit conversion; this chapter needs two.
The second is \term{widening} (\rr{widening}): an \code|int| is
converted into a \code|real| where a \code|real| is required. Both
conversions can apply one after the other: in \code|r := n| below, the
variable \code|n| is first dereferenced to an \code|int|, which is then
widened to a \code|real|:

\example{ch02/conversions}

There is no implicit conversion in the other direction, from \code|real| to
\code|int|, since it would lose information:

\example{ch02/no-narrowing}

The message, again in the terms of the Report, says that a \code|real|
cannot be converted to \code|int| in a ``strong-unit''. A \rr{unit} is an
expression; \rr{strong} describes the place where the expression appears.

This is the last important idea of the chapter: which implicit
conversions are allowed depends on where an expression stands. Algol 68
distinguishes five kinds of place, called \term{conversion contexts}
(\rr{syntactic positions}, \rr{contexts}). The right-hand side of a
declaration or an assignment is a \term{strong} context, where all
conversions are allowed, since the required type is known. An operand of
an operator is a \term{firm} context, where widening is \emph{not}
allowed, since the operator to apply is chosen by the types of its
operands and is not known in advance. The operator \code|ROUND|, for
example, is defined for a \code|real| operand only, and an \code|int|
operand is not widened to make it fit:

\example{ch02/round-int}

How then can \code|n + 0.5| work, with an \code|int| and a \code|real|
operand? Not by widening: the standard library simply defines the
arithmetic and comparison operators for each combination of \code|int| and
\code|real| operands (Revised Report, Section 10.2.3.5). Chapter~\ref{ch-advanced} describes
all six implicit conversions and all five contexts.

\section{Assignment}\label{sec-assignment}

An assignment \code|x := e| evaluates \code|e| and stores its value in
the cell that the reference \code|x| refers to. An assignment is itself
an expression: its value is the reference on the left-hand side. That
makes \code|a := b := 7| possible: it assigns 7 to \code|b|, and the result,
the reference \code|b|, is dereferenced and its value assigned to
\code|a|. The two sides of an assignment are evaluated in an unspecified
order (Section~\ref{sec-collateral}).

The \term{compound assignment} \code|x +:= e| adds \code|e| to the value
of \code|x| and stores the result in \code|x|; it is short for
\code|x := x + e|, except that \code|x| is evaluated only once. Similar
operators exist for the other arithmetic operators: \code|-:=|,
\code|*:=|, \code|/:=| (for a \code|real| variable), \code|%:=| and
\code|%*:=| (the \code|MOD| form). They are operators of precedence 1, so
\code|x +:= 1 + 2| adds 3. Their value, like that of an assignment, is
the reference on the left-hand side.

Unlike the assignment, these are ordinary operators, and like all
operators of the same precedence they are applied from left to right:

\example{ch02/plusab}

\noindent prints 6 and 1: \code|a +:= b +:= 3| is \code|(a +:= b) +:= 3|,
and \code|b| is not changed. A chain of assignments \code|a := b := 3|,
in contrast, groups to the right. And \code|n +:= x| for an \code|int|
variable \code|n| and a \code|real| \code|x| is rejected, for the same
reason as \code|n := n + x|.

\example{ch02/assignment}

\section{Explicit conversions and other operators}\label{sec-explicit}

A \code|real| is converted into an \code|int| with one of two operators.
\code|ENTIER x| is the largest integer not greater than \code|x| (the
floor). \code|ROUND x| is the integer nearest to \code|x|. Characters are
converted to and from their integer codes with \code|ABS| and
\code|REPR|. A few more operators complete the set used in this chapter:

\begin{center}
\begin{tabular}{ll}
\toprule
Operator & Result \\
\midrule
\code|ENTIER x| & the largest \code|int| not greater than the \code|real| \code|x| \\
\code|ROUND x| & the \code|int| nearest to the \code|real| \code|x| \\
\code|ABS x| & the absolute value of an \code|int| or \code|real| \\
\code|SIGN x| & $-1$, 0 or 1, according to the sign of \code|x| \\
\code|ODD n| & \code|true| if the \code|int| \code|n| is odd \\
\code|ABS c| & the character code of the \code|char| \code|c| \\
\code|REPR n| & the character with the code \code|n| \\
\bottomrule
\end{tabular}
\end{center}

\example{ch02/explicit}

\begin{implnote}
The Report does not say whether a number exactly halfway between two
integers, such as 2.5, is rounded up or down. ga68 and a68j round it away
from zero: \code|ROUND 2.5| is 3 and \code|ROUND -2.5| is $-3$.
\end{implnote}

\section{Comments and layout}\label{sec-comments}

In the convention of this book, a \term{comment} is text between
braces, \code|{ like this }|. Comments can be nested: \code|{ a { nested } comment }| is one comment. The
layout of a program, its spaces and line breaks, is free: a program can be
written on one line or spread over many. In the stropping of this book,
spaces are not allowed inside an identifier or a reserved word, but may
appear between the digits of a number. The Report allows spaces inside
identifiers as well, which is why it and older texts write
\code|max int| where this book writes \code|max_int|.

\begin{implnote}
The Report encloses a comment in two \code|#| symbols, or between two
of the bold words \textbf{comment} or \textbf{co}. In the stropping
convention of this book, ga68 and a68j accept only braces. Algol 68 Genie
accepts \code|#|\ldots\code|#|, \texttt{CO}\ldots\texttt{CO} and
\texttt{COMMENT}\ldots\texttt{COMMENT}, but not braces. These comments
cannot be nested: a comment ends at the next symbol of the kind that
began it.
\end{implnote}

\section{Simple input and output}\label{sec-io}

The procedure \code|print| writes a list of values to the standard
output. Each value is written in a standard layout: an \code|int| with
its sign, in a field wide enough for \code|max_int| (11 characters in ga68
and a68j), a \code|real| in scientific notation,
a \code|bool| as \code|T| or \code|F|, and a \code|char| or a string as
it is. Apart from the padding of an \code|int|, no separators are
inserted; a program that wants spaces between numbers prints them. For more control, three procedures convert a number
into a string:
\begin{itemize}
\item \code|whole (n, w)| converts the \code|int| \code|n| into a string
  of width \code|w|. A width of 0 gives a string just as wide as needed; a
  positive width gives a string of exactly that width, with a sign; a
  negative width gives that width without a plus sign.
\item \code|fixed (x, w, d)| converts the number \code|x| into a string of
  width \code|w| with \code|d| digits after the decimal point. Again,
  width 0 means ``as wide as needed''.
\item \code|float (x, w, d, e)| converts \code|x| into scientific
  notation with an exponent of \code|e| digits (Chapter~\ref{ch-transput}).
\end{itemize}

\example{ch02/output}

A number that does not fit into the given width is not printed in a
wider field: \code|whole (12345, 3)| yields three asterisks
(Section~\ref{sec-conversions}).

\begin{implnote}
The \code|print| of the Report writes a blank before every number that
does not begin a line. ga68, a68j and Algol 68 Genie leave it out, which
is why the outputs in this book show numbers directly after a preceding
string.
\end{implnote}

The procedure \code|read| reads values from the standard input into
variables. Numbers are separated by blanks or line breaks. The procedures
\code|read_int| and \code|read_real|, an extension of Algol 68 Genie that
\code|Genie_Transput| also provides, read a single value and return it,
which is convenient for initializing a constant:

\program{ch02/input}

With the input
\lstinputlisting[style=output]{examples/ch02/input.in}
\noindent the program writes
\outputof{ch02/input}

A number is read after skipping blanks and line breaks, but a
\code|char| is read as it is: a blank is a character. At the end of a
line, reading a \code|char| moves on to the first character of the next
line; the line break itself is never read as a character. The layout
procedure \code|space| skips a single character, so a
date written as \texttt{7/3/1975} can be read into three numbers:

\example{ch02/read-date}

\noindent prints \code|1975-3-7|.

\section{A complete program}\label{sec-complete}

The last program of this chapter reads a temperature in degrees
Fahrenheit and converts it to degrees Celsius. The output shown is for the
input \texttt{98.6}:

\example{ch02/temperature}

Writing \code|(fahrenheit - 32) * (5 / 9)| would also have worked, since
\code|/| yields a \code|real| even for two \code|int| operands; with
integer division, \code|5 % 9| would have been 0.

\section*{Exercises}
\addcontentsline{toc}{section}{Exercises}

The \hyperref[ans-ch2]{answers} are at the end of the book.

\begin{enumerate}
\item Which of the following are \code|int| literals, which are
  \code|real| literals, and which are neither?
  \code|12|, \code|-3|, \code|1 000|, \code|1,000|, \code|2.|, \code|.5|,
  \code|3e2|, \code|0.0|.
\item What are the values of \code|17 % 5|, \code|-17 % 5|,
  \code|17 MOD 5|, \code|-17 MOD 5| and \code|17 / 5|?
\item What are the values of \code|2 + 3 * 4 ** 2| and
  \code|-3 ** 2|?
\item Explain the difference between \code|int n = 10| and
  \code|int n := 10|. What is the type of \code|n| in each case?
\item Write the declaration \code|real x := 1.5, y| without the short
  form, using \code|loc|.
\item Which of the following are correct, given \code|int i := 1| and
  \code|real r := 2.5|? \quad (a) \code|r := i| \quad (b) \code|i := r|
  \quad (c) \code|i := ROUND r| \quad (d) \code|r := ROUND i|
\item After \code|int a := 1, b := 2|, what are the values of \code|a|
  and \code|b| after \code|a := b := a + b|?
\item Why does the compiler reject \code|maxCount := 0|, even if a
  variable \code|max_count| has been declared?
\item What does \code|int x := 1; ref int r = x; r := r + 1; print (x)|
  print?
\item Add parentheses to show how
  \code|a < b AND x > 0 = y > 0 OR NOT done| is grouped.
\item What are the values of \code|2.0 ** -2| and \code|ABS (-7 MOD 3)|?
  Why is \code|NOT n = 0| rejected for an \code|int| \code|n|?
\end{enumerate}

\section*{Programming problems}
\addcontentsline{toc}{section}{Programming problems}

\begin{enumerate}
\item Read a number of seconds and print it as hours, minutes and
  seconds.
\item Read the three coefficients of a quadratic equation whose roots
  are real, and print the roots (\code|sqrt| computes a square root).
\item Read a character and print its code, the next character, and
  whether it is a digit (\code|"0" <= c AND c <= "9"|).
\item Read an amount of money in cents and print the smallest number of
  coins of 50, 20, 10, 5, 2 and 1 cents that make it up.
\item \textbf{Easter.} Read a year and print the date of Easter Sunday,
  computed with the anonymous Gregorian algorithm (Meeus/Jones/Butcher),
  which needs only \texttt{\%} and \code|MOD|.
\item \textbf{Day of the week.} Read a date written as
  \texttt{dd/mm/yyyy} and print the day of the week as a number, using
  Zeller's congruence. Where does it matter that \code|MOD| never yields
  a negative result?
\item \textbf{Clock.} Read an hour from 0 to 23 and print it on a
  12-hour clock with \texttt{AM} or \texttt{PM}, without \code|if|: the
  hour is \code|(h + 11) MOD 12 + 1|. (Hint: \code|REPR| and
  \code|ABS| can compute the letter.)
\item \textbf{Rounding money.} Read an amount as a \code|real|, round it
  to cents with \code|ROUND|, and print euros and cents with \texttt{\%}
  and \code|MOD|. Try a negative amount, and an amount above 100000
  with the 32-bit \code|real| of ga68 or a68j. What goes wrong, and why
  should money not be kept in a \code|real|?
\end{enumerate}
